\section{The raw Euler-conjugated cutoff is unbounded on a narrow line}
\label{sec:euler-cutoff-obstruction}

Unboundedness of the canonical inverse in MA-MO-001 does not
by itself prove unboundedness inside a cutoff conjugation.
We now prove that stronger assertion for the actual one-sided
Euler multiplier model. The argument uses nonzero universal
values, not an assumed zero on the line.

Fix \(1/2<\sigma<1\), put \(a=\sigma-1/2\), and use
\[
 \mathcal H_a=L^2(\mathbb R,e^{2ax}\,dx),\qquad
 V_a=\mathcal F_+M_{e^{ax}},\qquad
 \mathcal F_+f(t)=(2\pi)^{-1/2}\int e^{itx}f(x)\,dx.
 \tag{MA-MO.20}
\]
Let \(L_0(s)\) be zeta or a fixed primitive Dirichlet
\(L\)-function. For a finite set \(S\) of finite primes put
\[
 \begin{aligned}
 A_0(t)&=L_0(\sigma+it),&
 E_S(t)&=\prod_{p\in S}(1-\chi(p)p^{-\sigma-it}),\\
 A_S(t)&=E_S(t)A_0(t),&
 j_S(t)&=A_S(t)^{-1}.
 \end{aligned}
 \tag{MA-MO.21}
\]
For zeta, \(\chi(p)=1\); a ramified Dirichlet value zero
gives the factor one. The closed operators \(Z_{a,S}\)
and \(J_{a,S}\) are the maximal multipliers \(M_{A_S}\)
and \(M_{j_S}\), conjugated by \(V_a\), as in MA-RS.9
and its character version MA-RS-005.

Take any nonconstant bounded \(\phi\in C^\infty(\mathbb R)\)
with compactly supported derivative. Its limiting constants
may agree; nonconstant compact position cutoffs are allowed.
Put
\[
 T_\phi=\mathcal F_+M_\phi\mathcal F_+^{-1},\qquad
 p(v)=\int e^{ivx}\phi'(x)\,dx,\qquad
 q(v)=\frac{i\,p(v)}{2\pi v}\quad(v\ne0),
 \tag{MA-MO.22}
\]
and define the dense Fourier core
\[
 \mathscr C_\sigma=C_c^\infty
    \bigl(\mathbb R\setminus\{t:L_0(\sigma+it)=0\}\bigr).
 \tag{MA-MO.23}
\]
No assertion that this line is zero free is needed.

\paragraph{MA-MO-002 (no bounded raw cutoff conjugation).}
The actual composition
\[
                    \mathcal Q_{\phi,S}
                         =Z_{a,S}M_\phi J_{a,S}
 \tag{MA-MO.24}
\]
is well defined on \(V_a^{-1}\mathscr C_\sigma\).
Neither it nor \(M_\phi-\mathcal Q_{\phi,S}\) has a
bounded extension to \(\mathcal H_a\).
In the Fourier picture there are unit vectors
\(v_M\in\mathscr C_\sigma\) with
\[
 \|M_{A_0}T_\phi M_{A_0^{-1}}v_M\|_2
                         \ge c\,e^{Md}/M\longrightarrow\infty
 \tag{MA-MO.25}
\]
for fixed \(c,d>0\). Finite Euler deletion changes only
a fixed positive comparison constant. The result remains
true on a line which is known to have no zero.

\begin{proof}
\emph{Domains before kernels.}
The line-zero set is discrete and has measure zero.
Removing small neighborhoods of the finitely many zeros
in each compact interval proves density of
\(\mathscr C_\sigma\) in \(L^2\).
The finite Euler factors do not vanish when \(\sigma>0\).
For \(v\in\mathscr C_\sigma\), \(j_Sv\in C_c^\infty\).
All derivatives of \(\phi\) are bounded, so \(M_\phi\)
preserves Schwartz space and \(T_\phi(j_Sv)\) is Schwartz.
The unconditional fixed-line polynomial growth of the
fixed zeta or Dirichlet function, already used in
MA-GW-001/004 and MA-RS, makes
\(A_ST_\phi(j_Sv)\) square integrable.
This proves every domain inclusion needed for (MA-MO.24).
No domain assertion for \(J_{a,S}M_\phi\) is being
silently inferred.

\emph{The actual cross-window kernel.}
The distribution kernel of \(T_\phi\) is
\((2\pi)^{-1}\widehat\phi_+(t-u)\).
Since \(\widehat{\phi'}_+(v)=-iv\widehat\phi_+(v)\),
the kernel off the diagonal is \(q(t-u)\).
On disjoint compact input and output windows the
already defined composition for \(S=\varnothing\)
therefore has the ordinary kernel
\[
                      \frac{A_0(t)}{A_0(u)}q(t-u).
 \tag{MA-MO.26}
\]
There is no delta or principal-value term on these
disjoint windows. Smooth core pairings can be computed
from this kernel.

\emph{The exact universality statement used.}
Voronin (1975), the unnumbered theorem in \S I,
printed pp.~443--444, gives uniform approximation of
every nonvanishing analytic target on \(|z|\le r<1/4\)
by \(z\mapsto\zeta(z+3/4+iT)\) for a real \(T\).
The following Remark on p.~444 gives the same assertion
for Dirichlet \(L\)-functions.
We use precisely those printed forms. No density
theorem or assertion that the approximating shifts
can be chosen arbitrarily large is required.
The classical universality proofs are published
background; the operator deduction is proved here.

Set \(h=\sigma-3/4\), and choose
\[
                      |h|<r<1/4 .
 \tag{MA-MO.27}
\]
The transform \(p\) is entire because \(\phi'\) has
compact support. It is not identically zero, by Fourier
injectivity and nonconstancy of \(\phi\).
The identity theorem therefore permits a fixed
sufficiently small \(d>0\) with
\[
                         h^2+d^2<r^2,\qquad q(d)\ne0.
 \tag{MA-MO.28}
\]
This also handles \(p(0)=0\) when the limiting values
of the cutoff agree.

Fix \(0<\epsilon<1/100\). For each positive integer \(M\),
the entire nonvanishing target
\[
                         f_M(z)=e^{-iM(z-h)}
 \tag{MA-MO.29}
\]
has minimum modulus \(e^{-Mr}\) on the disk.
Apply the printed theorem with absolute tolerance
\(\epsilon e^{-Mr}\). It gives a real \(T_M\) such that
\[
 L_0(z+3/4+iT_M)=f_M(z)(1+\eta_M(z)),\qquad
        |\eta_M(z)|<\epsilon\quad(|z|\le r).
 \tag{MA-MO.30}
\]
The relative error follows by division by the
displayed minimum modulus; it is not a stronger
universality input. On \(z=h+iv\) the target is \(e^{Mv}\).
In particular the windows below contain no line zero.

\emph{A lower bound for the actual operator norm.}
Choose a fixed nonnegative real
\(b\in C_c^\infty((-1,1))\), with \(\|b\|_2=1\), and put
\[
 \begin{aligned}
 v_M(u)&=M^{1/2}b(M(u-T_M)),\\
 w_M(t)&=M^{1/2}b(M(t-T_M-d)).
 \end{aligned}
 \tag{MA-MO.31}
\]
These are unit vectors. For large \(M\) their supports
are disjoint and lie in the two vertical windows strictly
inside the fixed disk. In particular \(v_M\) is in the
declared core.
Substitute \(u=T_M+\xi/M\) and
\(t=T_M+d+\eta/M\) in (MA-MO.26). With the inner
product linear in its first argument, the pairing is
\[
 \begin{split}
 &\langle M_{A_0}T_\phi M_{A_0^{-1}}v_M,w_M\rangle\\
 &\quad=\frac{e^{Md}}{M}\iint b(\xi)b(\eta)e^{\eta-\xi}
       q\left(d+\frac{\eta-\xi}{M}\right)
       \frac{1+\eta_M(h+i(d+\eta/M))}
            {1+\eta_M(h+i\xi/M)}\,d\xi\,d\eta .
 \end{split}
 \tag{MA-MO.32}
\]
The error ratio differs from one by at most
\(2\epsilon/(1-\epsilon)\).
Continuity of \(q\) and \(q(d)\ne0\) give, for large \(M\),
\[
 \sup_{|\xi|,|\eta|\le1}
 \left|q\left(d+\frac{\eta-\xi}{M}\right)/q(d)-1\right|<1/100.
\]
The product of these two normalized factors is within
\(1/2\) of one. The remaining weight is nonnegative,
with strictly positive integral
\[
 I_b=\left(\int b(\eta)e^\eta d\eta\right)
                   \left(\int b(\xi)e^{-\xi}d\xi\right)>0.
 \tag{MA-MO.33}
\]
Subtracting the main value \(q(d)I_b\) inside the
double integral proves
\[
 |\langle M_{A_0}T_\phi M_{A_0^{-1}}v_M,w_M\rangle|
               \ge\tfrac12|q(d)|I_b\,e^{Md}/M.
 \tag{MA-MO.34}
\]
The unit norm of \(w_M\) proves (MA-MO.25).
This argument proves a norm lower bound on actual
core vectors; it does not infer it from a large
kernel value at a single point.

\emph{Finite Euler factors.}
There are constants \(0<m_S\le M_S<\infty\) such that
\(m_S\le|E_S(t)|\le M_S\) for all real \(t\).
Use the factors \(1-|\chi(p)|p^{-\sigma}>0\)
and \(1+|\chi(p)|p^{-\sigma}\).
Both \(E_S\) and its inverse preserve the core, and
\[
 M_{A_S}T_\phi M_{j_S}
       =M_{E_S}
          (M_{A_0}T_\phi M_{A_0^{-1}})
                                M_{E_S}^{-1}.
 \tag{MA-MO.35}
\]
This is an exact bounded similarity on that core.
The normalized inputs \(E_Sv_M/\|E_Sv_M\|_2\)
give at least \(m_S/M_S\) times the lower bound
(MA-MO.25). Thus no finite \(S\) has a bounded
extension. Since \(T_\phi\) is bounded, subtracting
it cannot turn this unbounded operator into a
bounded one. Conjugation by \(V_a\) finishes the proof.
\end{proof}

\paragraph{MA-MO-003 (the rational smooth-image graph transfer).}
In the rational sector the obstruction already holds
on actual smooth theta-image vectors. More precisely,
the graph closure of the restriction of
\(\mathcal Q_{\phi,S}\) to those vectors on which its
closed multiplier composition is defined contains
its graph on \(V_a^{-1}\mathscr C_\sigma\).
Consequently no bounded estimate on that smooth image
can extend the declared composition to \(\mathcal H_a\).
This remains a statement about the specified multiplier
continuation of its output.

\begin{proof}
Fix \(v\in\mathscr C_\sigma\) and let \(q_0=j_Sv\),
a smooth compact function. Choose \(h_0\in\mathscr S(\mathbb R)\)
with
\[
          \mathcal F_+h_0(t)=q_0(t)/(1-\sigma-it).
 \tag{MA-MO.36}
\]
There is no real zero of this denominator.
For \(\epsilon>0\) put
\[
 H_\epsilon(x)=e^{-\epsilon x^2-ax}h_0(x),\qquad
 F_\epsilon=(\partial_x+\tfrac12)H_\epsilon .
 \tag{MA-MO.37}
\]
These functions belong to \(\mathcal E\).
Integration by parts gives
\(\mathcal M F_\epsilon(1/2)=0\).
Moreover
\[
 q_\epsilon:=V_aF_\epsilon
       =(1-\sigma-it)\mathcal F_+(e^{-\epsilon x^2}h_0)
                           \longrightarrow q_0
                               \quad\hbox{in Schwartz space}.
 \tag{MA-MO.38}
\]
For completeness, Gaussian multiplication tends to the
identity in every Schwartz seminorm: split into a fixed
compact interval and its complement; derivatives converge
uniformly on the compact part, while the product rule,
the Schwartz tails of \(h_0\), and uniformly polynomially
bounded derivatives of \(e^{-\epsilon x^2}\) for
\(0<\epsilon\le1\) control the complement.
Fourier transformation and the linear multiplier preserve
this convergence.

To see the actual theta-image assertion, define the even
archimedean function, first for \(t>0\), by
\[
                   f_\epsilon(t)=
                    F_\epsilon(\log t)/(2\sqrt t).
 \tag{MA-MO.39}
\]
Superexponential decay with all derivatives makes it
Schwartz and flat at zero. Its full real integral is
\(\mathcal M F_\epsilon(1/2)=0\).
Thus both theta polar conditions hold, and the actual
theta formula of MA-TR-101/MA-GP.1 gives
\[
 (\Theta f_\epsilon)(x)
       =\sum_{n\ge1}n^{-1/2}F_\epsilon(x+\log n)
                       \in T_{1/2}\mathcal I .
 \tag{MA-MO.40}
\]
On a right Mellin line its transform is the product
of zeta and the transform of \(F_\epsilon\).
The moment condition removes the sole zeta pole.
The actual smooth theta identity therefore continues
to the line \(\sigma\).
Finite Euler deletion is a finite combination of
translations and preserves this smooth ideal.
Consequently the actual smooth-image vector
\(g_\epsilon=Z_{a,S}F_\epsilon\) has
\(V_ag_\epsilon=A_Sq_\epsilon\).
It is in \(\mathcal D(J_{a,S})\), and
\(J_{a,S}g_\epsilon=F_\epsilon\), since multiplication
by \(j_S A_S\) is one almost everywhere.

The convergence (MA-MO.38), polynomial growth of \(A_S\),
and continuity of \(T_\phi\) on Schwartz space now give
\[
 \begin{aligned}
 V_ag_\epsilon&=A_Sq_\epsilon\longrightarrow A_Sq_0=v
                                                   &&\text{in }L^2,\\
 V_a\mathcal Q_{\phi,S}g_\epsilon
       &=A_ST_\phi q_\epsilon
                        \longrightarrow A_ST_\phi q_0
                                                   &&\text{in }L^2.
 \end{aligned}
 \tag{MA-MO.41}
\]
In particular each \(g_\epsilon\) lies in the composition
domain and its output lies in \(\mathcal H_a\), regardless
of whether that output is in the smooth ideal.
This is exactly graph approximation.
A bounded estimate on these smooth-image vectors would
pass to every core vector in (MA-MO.41), contradicting
MA-MO-002. The same conclusion holds for the difference
with \(M_\phi\), since that multiplier is bounded.
\end{proof}

\paragraph{The exact common-space consequence and its limits.}
For Meyer's source cutoff, the first diagonal entry of
the centered difference \(P_+^S-P_-\) is
\[
 D=Z_SM_{1-\phi}J_S-M_{1-\phi}
                   =M_\phi-Z_SM_\phi J_S
 \tag{MA-MO.42}
\]
on its source domain, as derived from the definitions
in MA-MEY.13--.14. The exact locator is Meyer (2005),
Lemmas~5.3--5.4 and the following projection displays,
then \texttt{eq:global\_perturbation} in the proof of
Lemma~5.6, author-manuscript pp.~39--42.
The already recorded sign correction is retained.

A bounded operator on
\(\mathcal H_a\oplus\mathcal H_{-a}\), \(0<a<1/2\),
which agrees with this multiplier continuation on its
declared core, would give a bounded \(D\) by coordinate
compression. MA-MO-002 rules that out. MA-MO-003 also
rules out obtaining the bound merely by restricting
this composition to the rational smooth theta image.
This settles that specific bounded narrow matrix route.
It does not identify the multiplier-continued output,
with its possible pole subtraction, with every
uncontinued geometric inverse composition.
The subsequent MA-MO-004 identifies the exact rational source
polar correction, and MA-MO-005 removes it on actual smooth
theta-image witnesses with unbounded literal raw outputs.
Those strengthened statements retain the distinction between
the full direct product and the smaller common-sum subspace.

Meyer's wide-weight nuclear theorem uses \(a>1/2\)
after centering and is unaffected.
On the generic lines and Gevrey cutoffs where MA-RS
checks the literal commutator domains,
\(D=Z_{a,S}[J_{a,S},M_\phi]\) on the stated core;
the present result proves this unsmoothed operator
unbounded. Its product with \(C_g\) can still be
trace class, as MA-RS-002 proves.
No unclosability assertion or impossibility of
every nonlocal transport is made.
The one-sided \(\mathcal H_a\) is not substituted
for the two-sided weighted PW space or its Gram factors.
The separate MA-MO-009 supplies the required primary
Hecke universality input and proves the non-induced
imaginary-quadratic extension, including literal
smooth ideal-sum inputs and outputs.

\paragraph{Trace-framework placement and source boundary.}
This is an analytic obstruction to bounded entries
of the specified Hilbert matrix, not a new algebraic
obstruction to the finite-length modules or their
\(K_0\) and exterior-power identities.
In the inherited cote~1 nuclear-trace framework,
scans~079--080 and~111--112 through
ALPOGE\_BRIDGES BR-TR-01/02, the relevant topology
must be justified before cyclicity is used.
The smoothed constructions do so on their stated
domains; the raw narrow matrix has no such bounded
extension. No motivic positivity or weight theorem
is inferred, and no zero-location theorem is an input.

Voronin, S. M. (1975). Theorem on the universality
of the Riemann zeta-function.
\emph{Mathematics of the USSR-Izvestiya, 9}(3), 443--453.
\url{https://doi.org/10.1070/IM1975v009n03ABEH001485}.
The exact disk theorem in \S I, pp.~443--444, and
the Dirichlet Remark on p.~444 were freshly read
in retrieved primary text. This is a statement-scope
check; the full classical universality proof remains
published background. The new deductions above
receive separate whole-proof checks.
